% appendix_a_reproduction.tex -- Appendix A Reproduction: model names, environment, runs and job ids, scripts,
% commands. Owner: Appendix writer. The ONLY chapter where job ids and old run names appear (REWRITE_PLAN.md section 5).
% Sources: docs/handover.md 0a (commands of 18 Oct, NAVER baseline, model cost, student), 0b, 1, 3.6, 4f, 8 (job
% index), 10; docs/deck_explained_2026-10-19.md slide 30; data/CLUSTER_ANSWERS.md; the script headers named below;
% slurm/env.sh; slurm/egl.sh; exp7_final_model.sh :74-92 (GPUs), :853 (container); exp8_distill_vitb.sh header.
% Tables: T19 (tables/T19_model_names.tex, first), T13 (tables/T13_runs.tex). docs/model_names.md = T19 in Markdown.
% Trimmed 19 Oct: the scripts table (old A.4, tab:scripts-new) and the commands (old A.5) moved to docs/thesis/REPRODUCE.md;
% labels sec:app-scripts and sec:app-commands now sit on the pointer at the end of the runs section. Long version:
% git show 76aa1a4:docs/thesis/chapters/appendix_a_reproduction.tex.
% 12 Nov revision (REVISION_BRIEF_2026-11-12.md): opening states that the main text uses descriptive names; name key
% moved here from the T19 caption, path key from the T13 caption; sweep status (CA "8 Nov"), ablation arms and the
% student's Relative Human run added as pending; first person removed; run-to-run variation by reference to
% sec:method-variation.
% 6 Nov additions: exp9 header; exp8_distill_vitb.sh:257, :932, :1142 (seed only in the bone-fix trace; train.py
% has no manual_seed); docs/video_inference.md:36-60, :387-440; demo_video.py:826-843, :911-926. Pending: F6-F8.

\chapter{Model names, runs and reproduction}
\label{app:reproduction}
\begingroup % until the end of the file: \path may also break after _ and - (long run names)
\expandafter\def\expandafter\UrlBreaks\expandafter{\UrlBreaks\do\_\do\-}%

The main text refers to every model by a descriptive name, such as the teacher, the student or the initial model.
This appendix collects the run identifiers behind these names, together with the Slurm jobs, logs, checkpoints and
files that produced each number.

Every number can be traced in two steps. Each headline number is a macro in one of the files \path{numbers_*.tex} of
the thesis source, and the comment next to it names the job and the line of its log, for example
\texttt{3611336:191}. A job is one unit of work run by Slurm, the scheduler of the cluster, and the job id is the
number Slurm assigned to it. Table~\ref{tab:names} maps every descriptive name to its run identifier and to the run
name recorded during the project; Table~\ref{tab:runs} maps each job to the model, the log and the checkpoint. Job
ids and recorded run names appear only in this appendix.

\section{Model names}
\label{sec:app-names}

The ViT-L line (``L'') comprises every model that keeps the ViT-L backbone. L-stage 1 to 4 are the steps of the
transfer schedule that turned the released Multi-HMR-Anny checkpoint (the released baseline) into the initial model (Base), which is the end of that
schedule and the start of every later fine-tune. The teacher (Final) and the teacher without corrections (Control)
are the two arms of the last teacher experiment (exp7): the same \MethodFinalEpochs-epoch fine-tune of Base, with and
without the two corrections. The number after ``Student-S/'' is the old version number of a ViT-S student; the
student (Student) is the ViT-B model distilled from the teacher. The 60-epoch student (Student-60) is the student's
training configuration trained twice as long (Section~\ref{sec:res-training-length}); it is not the reported model.
The sweep of Section~\ref{sec:res-kd-temperature} trains the same configuration at \KDSweepArms{} temperatures.
The two arms of the distillation ablation (Section~\ref{sec:res-kd-ablation}), the student without distillation
losses (Student-noKD) and the ViT-B model without teacher (ViT-B-noTeacher), are the last rows of the table.
Two evaluation protocols are used (Section~\ref{sec:method-eval}): the holdout protocol, with the helper bones set to
rest pose before decoding, and the 3DPW protocol, with the helper bones scored as predicted. File and log names use V1
for the holdout protocol and V0 for the 3DPW protocol.

In Table~\ref{tab:names}, ``ep.''\ is the epoch kept. Backbone: L, B, S = DINOv2 ViT-L, ViT-B, ViT-S with 14-pixel
patches; ``all'' = every block trained, ``last 4'' = the last 4 of 24 blocks, ``+p'' = a linear projector to the
heads' width 1024. Data: A = AnnyOne; A+P = 70\,\% AnnyOne and 30\,\% 3DPW train, with the old labels (before their
correction, Chapter~\ref{ch:data}) or the corrected v23 labels. Schedule: epochs $\times$ iterations per epoch
$\times$ batch, then the learning rate, halved every 30 epochs unless noted (planned epochs where the run stopped
early). Checkpoints lie in \path{logs/anny_model/<run>/checkpoints/} with copies in \path{models/multiHMR/};
Table~\ref{tab:runs} gives the path of every model with a number in the thesis, and Appendix~\ref{app:flags} the
full flags of Final, Control and the student. ``Not recorded'' means that no script, log or document in the
repository states the value.

% T19 Model names: descriptive name of the main text -> run identifier -> old run / checkpoint name, start, backbone,
% data, schedule (and the one setting that is new in the row). Owner: Appendix writer. First table of Appendix A.
% 12 Nov revision (REVISION_BRIEF_2026-11-12.md section 2 rule 4): the main text uses descriptive names only; this
% table is the single place that maps them to the run identifiers (Student, Student-60, Final, Control, Base,
% Multi-HMR-Anny, Final-masked, Final-partial-mask, Mixed-L, Fresh-heads-L, L-stage 1-4, Student-S/*). First column =
% descriptive name, identifier in bold below it. Rows for the two distillation-ablation arms added (brief section 5:
% (a) distillation losses off, teacher heads still copied; (b) no teacher, no losses, no head copy; submission not confirmed,
% run names fixed by slurm/experiments/exp10_kd_ablation.sh "RUN NAMES"; not completed at submission, fix pass D1). Student-60 was called "Student, 60 epochs" before 12 Nov. The table key moved from
% the caption into the paragraph before the table (appendix_a_reproduction.tex, sec:app-names).
% 13 Nov: ablation rows completed (jobs 3744297, 3744307; CA "13 Nov: exp10 results" A); results in tab:kd-ablation.
% Same content as docs/model_names.md (for the deck and the handover).
% Trimmed 19 Oct to one longtable (was two: part 2 held the full flag lists and checkpoint paths; recover with
% git show 76aa1a4:docs/thesis/tables/T19_model_names.tex). Full flags of Final, Control and the Student: Appendix C;
% checkpoint paths of every model with a number in the thesis: T13.
% Sources:
%   Multi-HMR-Anny: slurm/history/train_step1_pretrained_bb.sh:12 (training data "annyone, bedlam, coco, mpii");
%     head size 76.57 M (job 3614947, H 0a); slurm/eval/run_eval_baseline_naver.sh header.
%   L-stage 1-4, Base: slurm/history/train_step1_pretrained_bb.sh (+ _resume, RESUME_EPOCH=177),
%     train_step2_frozen_bb.sh (+ _resume), train_s2_partialft.sh, train_s2_boosted_v4.sh (V3_EPOCH=00054),
%     train_s2_shape_v5.sh (V4_EPOCH=00299); README.md:236-246 (pipeline 2); H 1 "Exact checkpoint paths", "v5 line";
%     H 3.2 "Origin" (stage 2: fresh heads at 1e-4, resumed from its epoch 283 at 3e-5, saved args 3573519:3; the
%     history scripts were edited after the runs, so these recipes are as recorded there, not re-measured);
%     draft main.tex:565-569 and tab:main. Base lr 1.25e-6 in the last 10 epochs; L-stage 1 has no kept checkpoint.
%   Mixed-L: slurm/train/train_teacher_v9_mixed.sh (V5_EPOCH=00099, LABELS threedpw_anny_full/labels); H 1.
%   Fresh-heads-L: slurm/train/train_v10_scratch_mixed.sh:3-11 (despite "scratch": NAVER's DINOv2 ViT-L backbone and
%     fresh heads); H 1 table (epoch 290 traced, 3571794); copy anny_v10_scratch_mixed_00299.pt in
%     slurm/eval/run_eval_3dpw_test.sh MODELS list (draft tab:kd_eval: epoch 299).
%   Control, Final, Final-masked, Final-partial-mask: slurm/experiments/exp7_final_model.sh (header, :207-257, :759-771);
%     H 0a step 6; H 8 job index; jobs 3611339, 3611336, 3612325, 3614042 (done 19 Oct, CLUSTER_ANSWERS).
%   Student-S/1-6c, mixed, final: slurm/history/train_distill_vits{,_v2,_v3,_v4,_v5,_v6,_v6b,_v6c,_v7_mixed}.sh,
%     slurm/train/train_distill_vits_v8_final.sh; draft tab:student_runs, tab:kd_epochs, tab:kd_eval (best epochs);
%     checkpoint copies from slurm/eval/run_eval_3dpw_test.sh MODELS list. --brightness/--contrast/--saturation/--hue
%     are parsed by train.py (:1722-1725) and read nowhere (grep args.brightness: no hit), so not listed as flags.
%     Student-S/final was evaluated as evaltest_student_v8_final.
%   Student: slurm/experiments/exp8_distill_vitb.sh header and :319-331; job 3614191 (done 19 Oct; V1 holdout PVE 119.6
%     3614191:195-207, 3DPW test v23 PVE 110.1, final table 3614191:259-283 FAIL; data/CLUSTER_ANSWERS.md last section).
%     Typed 119.6 / 110.1 replaced by \StudentPVE / \StudentThreeDPW on 6 Nov.
%   Student-60 (added 6 Nov): job 3623342, RUN ..._e60, checkpoint 00059.pt; data/CLUSTER_ANSWERS.md "6 Nov: 60-epoch
%     student final table" and the paragraph after it (LR halved every 24 = exp8's default 2/5 of EPOCHS,
%     slurm/experiments/exp8_distill_vitb.sh:229; EVAL_FREQ 5; FAIL on all four bars); same starting weights and flags as the Student.
%   Temperature sweep (added 6 Nov): slurm/experiments/exp9_kd_temperature.sh header (TEMPS 1 2 4 8 16, RUN names
%     <prefix>_e30_kdT<t>_sweep in results/experiments/exp8_distill_vitb/, EVAL_FREQ 1, VAL_LOSS 1, everything else
%     exp8's defaults); jobs 3717507/3717530/3717538/3717581/3717598 (CA "8 Nov: exp9 temperature sweep done"; T=16
%     trained to the end, post-training step failed on an HTTP 504 time-out).
\begingroup
\expandafter\def\expandafter\UrlBreaks\expandafter{\UrlBreaks\do\_\do\-}% let \path break after _ and - here
\scriptsize
\setlength{\tabcolsep}{2.5pt}
\begin{longtable}{@{}>{\raggedright}p{1.95cm}>{\raggedright}p{2.75cm}>{\raggedright}p{2.5cm}>{\raggedright}p{0.95cm}>{\raggedright}p{1.15cm}>{\raggedright\arraybackslash}p{2.8cm}@{}}
\caption[Model names and run identifiers]{Descriptive name used in the main text, run identifier (bold) and
recorded run name of every model, with starting weights, backbone, training data and schedule. Key in
Section~\ref{sec:app-names}.}
\label{tab:names}\\
\toprule
Main-text name, \textbf{identifier} & Recorded run / checkpoint name & Starting weights & Backb. & Data & Schedule; what it is \\
\midrule
\endfirsthead
\toprule
Main-text name, \textbf{identifier} & Recorded run / checkpoint name & Starting weights & Backb. & Data & Schedule; what it is \\
\midrule
\endhead
\bottomrule
\endfoot
\multicolumn{6}{@{}l}{\emph{Released starting point and baseline}} \\
the released baseline\newline\textbf{Multi-HMR-Anny} & \path{multiHMR_672_L_anny} (released checkpoint) & trained by NAVER & L, own decoder & A, BEDLAM, COCO, MPII & not recorded; evaluated with its own model class \\
\addlinespace
\multicolumn{6}{@{}l}{\emph{ViT-L line (``L''): the transfer schedule and the initial model}} \\
earlier model\newline\textbf{L-stage 1} & \path{anny_s1_pretrained_bb_v2} & Multi-HMR-Anny backbone, fresh heads & L, all & A & 150 $\times$ 1,000 $\times$ 4, $10^{-4}$, from epoch 177 at $3\cdot10^{-5}$; side branch, no checkpoint kept \\
earlier model\newline\textbf{L-stage 2} & \path{anny_s2_frozen_bb_v2}, ep.~299 (``v2'') & Multi-HMR-Anny backbone, fresh heads & L, frozen & A & 150 $\times$ 1,000 $\times$ 4, $10^{-4}$, from epoch 283 at $3\cdot10^{-5}$ \\
earlier model\newline\textbf{L-stage 3} & \path{anny_s2_partialft_v3}, ep.~54 (``v3'') & L-stage 2, epoch 299 & L, last 4 & A & 300 $\times$ 1,000 $\times$ 4, $10^{-5}$; helper bones held at rest \\
earlier model\newline\textbf{L-stage 4} & \path{anny_s2_boosted_v4}, ep.~299 (``v4'') & L-stage 3, epoch 54 & L, last 4 & A & 300 $\times$ 1,000 $\times$ 4, $10^{-5}$; neck and hand bones weighted $\times$4 \\
the initial model\newline\textbf{Base} & \path{anny_s2_shape_v5}, ep.~99 (``v5'') & L-stage 4, epoch 299 & L, last 4 & A & 100 $\times$ 1,000 $\times$ 4, $10^{-5}$; shape weight 20; AnnyOne only, start of every later fine-tune \\
earlier model\newline\textbf{Mixed-L} & \path{anny_teacher_v9_mixed}, ep.~62 and 149 (``v9'') & Base & L, last 4 & A+P, old & 150 $\times$ 1,000 $\times$ 4, $10^{-5}$, halved every 60 \\
earlier model\newline\textbf{Fresh-heads-L} & \path{anny_v10_scratch_mixed}, ep.~290 (``v10'', not from scratch) & Multi-HMR-Anny backbone, fresh heads & L, last 4 & A+P, old & 300 $\times$ 1,000 $\times$ 4, $10^{-5}$; 3DPW from the first step, no neck or hand weight \\
\addlinespace
\multicolumn{6}{@{}l}{\emph{Teacher experiment (exp7): two arms and two shape-mask variants}} \\
the teacher without corrections (uncorrected teacher)\newline\textbf{Control} & exp7 \path{v5_root_dist1_nofix_v23_e10}, ep.~9; job 3611339 & Base & L, last 4 & A+P, v23 & 10 $\times$ 1,000 $\times$ 4, $10^{-5}$; 3DPW hands masked, pelvis targets \\
the teacher\newline\textbf{Final} & exp7 \path{v5_pc_dist5_bonefix_v23_e10}, ep.~9; job 3611336 & Base and its bone correction & L, last 4 & A+P, v23 & as Control, with head targets, depth and translation weights 5, bone correction \\
the teacher with full shape mask\newline\textbf{Final-masked} & exp7 \path{v5_pc_dist5_bonefix_v23_e10_shpmask}, ep.~9; job 3612325 & as Final & L, last 4 & A+P, v23 & as Final, 8 converter constants masked for 3DPW people \\
the teacher with partial shape mask\newline\textbf{Final-partial-mask} & exp7 \path{v5_pc_dist5_bonefix_v23_e10_shpmask_noage}, ep.~9; job 3614042 & as Final & L, last 4 & A+P, v23 & as Final, the 7 constants other than age masked; not adopted \\
\addlinespace
\multicolumn{6}{@{}l}{\emph{ViT-S students (``Student-S''), earlier models; teacher Base unless noted}} \\
earlier model\newline\textbf{Student-S/1} & distillation v1, \path{anny_distill_vits} & ImageNet DINOv2-S, random heads & S, all & A & 300 $\times$ 1,000 $\times$ 4, $10^{-4}$; feature distillation \\
earlier model\newline\textbf{Student-S/2} & v2, \path{anny_distill_vits_v2_a} & Student-S/1, epoch 199 & S, all & A & 200 $\times$ 1,000 $\times$ 4, $2\cdot10^{-5}$, halved every 60 \\
earlier model\newline\textbf{Student-S/3} & v3, \path{anny_distill_vits_v3} & ImageNet DINOv2-S; 54 of 67 head tensors from Base & S, all & A & 150 $\times$ 1,000 $\times$ 4, $5\cdot10^{-5}$, halved every 60; output distillation \\
earlier model\newline\textbf{Student-S/4} & v4, \path{anny_distill_vits_v4} & NAVER's \path{multiHMR_672_S} backbone; heads as /3 & S, all & A & 150 $\times$ 1,000 $\times$ 8, $7\cdot10^{-5}$, halved every 60 \\
earlier model\newline\textbf{Student-S/5} & v5, \path{anny_distill_vits_v5} & as Student-S/4 & S, all & A & as /4; heads frozen for 5 epochs \\
earlier model\newline\textbf{Student-S/6} & v6, \path{anny_distill_vits_v6}, ep.~137 & NAVER's ViT-S backbone; all 67 head tensors from Base & S+p & A & as /4; stopped at epoch 138 by the time limit \\
earlier model\newline\textbf{Student-S/6b} & v6b, \path{anny_distill_vits_v6b}, ep.~140 & Student-S/6, epoch 137 & S+p & A & 150 $\times$ 1,000 $\times$ 8, $2\cdot10^{-5}$, halved every 60; teacher in bfloat16 \\
earlier model\newline\textbf{Student-S/6c} & v6c, \path{anny_distill_vits_v6c}, ep.~149 & Student-S/6b, epoch 140 & S+p & A & 150 $\times$ 1,000 $\times$ 16, $3\cdot10^{-5}$, halved every 60 \\
earlier model\newline\textbf{Student-S/mixed} & v7, \path{anny_distill_vits_v7_mixed}, ep.~130 & Student-S/6b, epoch 140 & S+p & A+P, old & 150 $\times$ 1,000 $\times$ 8, $2\cdot10^{-5}$, halved every 60 \\
earlier model\newline\textbf{Student-S/final} & v8, \path{anny_distill_vits_v8_final}, ep.~52 & Student-S/6c, epoch 149; teacher Mixed-L & S+p & A+P, old & 150 $\times$ 1,000 $\times$ 16, $3\cdot10^{-5}$, halved every 60 \\
\addlinespace
\multicolumn{6}{@{}l}{\emph{ViT-B student and its sensitivity runs}} \\
the student\newline\textbf{Student} & exp8 \path{student_vitb_from_v5_pc_dist5_bonefix_v23_e10_00009_e30}; job 3614191 (done 19 Oct) & NAVER's \path{multiHMR_672_B} backbone; every head tensor from Final; teacher Final & B+p, all & A+P, v23 & 30 $\times$ 1,000 $\times$ 8, $3\cdot10^{-5}$, halved every 12; Final's targets and masks; holdout \StudentPVE, 3DPW \StudentThreeDPW\,mm, 95\,\% retention thresholds not met; the reported model \\
the 60-epoch student\newline\textbf{Student-60} & exp8 \path{student_vitb_from_v5_pc_dist5_bonefix_v23_e10_00009_e60}, ep.~59; job 3623342 & as Student & B+p, all & A+P, v23 & as Student, but \StudentSixtyEpochs{} epochs, halved every \LossLRHalfSixty; holdout \StudentSixtyPVE, 3DPW \StudentSixtyThreeDPW\,mm, 95\,\% retention thresholds not met; training-length analysis, not the reported model \\
the student at temperature $t$\newline\textbf{Student, $T=t$} & exp8 (submitted by exp9) \path{student_vitb_from_v5_pc_dist5_bonefix_v23_e10_00009_e30_kdT<t>_sweep}; jobs in Table~\ref{tab:runs} & as Student & B+p, all & A+P, v23 & as Student, with feature-distillation temperature $t \in \{1, 2, 4, 8, 16\}$; evaluation after every epoch and validation loss on (both only measure); temperature analysis \\
\addlinespace
\multicolumn{6}{@{}l}{\emph{Distillation ablation}} \\
the student without distillation losses\newline\textbf{Student-noKD} & exp8 \path{student_vitb_from_v5_pc_dist5_bonefix_v23_e10_00009_e30_nokd}, ep.~29; job 3744297 & as Student (teacher heads copied) & B+p, all & A+P, v23 & as Student, with both distillation losses off; holdout \AblNoKDHold, 3DPW \AblNoKDTP\,mm, 95\,\% retention thresholds not met; arm (a), results in Table~\ref{tab:kd-ablation} \\
the ViT-B model without teacher\newline\textbf{ViT-B-noTeacher} & exp8 \path{student_vitb_from_v5_pc_dist5_bonefix_v23_e10_00009_e30_supervised}, ep.~29; job 3744307 & NAVER's \path{multiHMR_672_B} backbone, fresh heads & B+p, all & A+P, v23 & as Student, with no teacher: no distillation losses, no head copy or freeze, no bone correction in training; holdout \AblNoTeachHold, 3DPW \AblNoTeachTP\,mm, 95\,\% retention thresholds not met; arm (b), results in Table~\ref{tab:kd-ablation} \\
\end{longtable}
\endgroup


\section{Environment}
\label{sec:app-env}

All training, evaluation and demo jobs ran on the DFKI Pegasus Slurm cluster, from the repository clone in
\path{/netscratch/najib/multi-hmr}. The repository is public on GitHub
(\path{saadnajib/Multi-Person-Whole-Body-Human-Mesh-Recovery-in-a-Single-Shot-with-Annyone-dataset-and-3dpw-dataset}).
Checkpoints, logs and datasets are not in it. Each job script prints the short git hash of the clone into the run's
\path{settings.txt}, so a run can be matched to a commit.

Every job runs inside the container image \path{ubuntu+20.04_v4.sqsh} (in \path{/netscratch/najib/ContainerImages/}),
started by \texttt{srun} with the option \path{--container-image} and the mounts \path{/fscratch}, \path{/netscratch},
\path{/home/najib} and \path{/dev/shm}. Inside it, \path{slurm/env.sh} is sourced first. It activates the conda
environment \texttt{multihmr} (Python 3.9, PyTorch 2.0.1 with CUDA 11.7, xformers 0.0.20, torchvision 0.15.2) and
moves the Hugging Face, Torch and Anny caches to \path{/netscratch/najib}. It also sets one thread per data-loader
worker and strips the \texttt{X | None} type hints from the downloaded DINOv2 hub code, because Python 3.9 rejects
them. \path{slurm/egl.sh} then installs the EGL libraries for the headless renderer. The video launcher also tests a
real render and stops before loading a checkpoint if that test fails.

The landmark detector of the video pipeline, MediaPipe 0.10.14, runs in its own Python 3.10 environment, with the
packages pinned in \path{requirements-video-landmarks.txt}. There, \path{scripts/video/extract_pose_landmarks.py} reads
the video once on the CPU and writes a JSON file before the GPU job starts. \path{docs/video_inference.md} creates this environment at
\path{/netscratch/najib/venvs/video-landmarks}.

That PyTorch build has GPU kernels up to compute capability 8.6. For this reason the training and evaluation scripts
ask for A100, A40, L40S and RTX A6000 nodes and refuse an H100 or H200 (compute capability 9.0). The V100 nodes
\texttt{gifu} and \texttt{glendale} are excluded: every converter task sent there failed within a minute with an
argument-parser error whose cause was not identified. Only the 3DPW converter ran on H100 nodes, through a small virtual
environment with PyTorch 2.1.2 for CUDA 12.1 on top of \texttt{multihmr}
(\texttt{PY\_VENV=/netscratch/najib/venv\_cu121}). Steps that need only NumPy (building the 3DPW label files, the re-fit
check) run on the login node. The model-cost script needs a GPU node, because \path{model.py} imports the renderer and
the login node has no EGL.

\section{Runs, jobs, logs and checkpoints}
\label{sec:app-runs}

Table~\ref{tab:runs} lists every job behind a number, a table or a figure of the thesis. Paths are relative to the
project folder on the cluster. The log of job $j$ is \path{logs/<log name>_}$j$\path{.out} (array jobs:
\path{_}$j$\path{_<task>.out}; the video and image jobs: \path{logs/<log name>-}$j$\path{.out}, with a hyphen).
\textsf{E7} stands for \path{results/experiments/exp7_final_model/}, \textsf{M} for \path{models/multiHMR/}.

Of the \KDSweepArms{} temperature-sweep jobs (submitted 6 November), four completed; the $T = \KDSweepMaxT$ job
(3717598) trained to the end, but its post-training evaluation failed on a network time-out, so it has no final table.
The per-temperature validation values of the sweep were logged by the runs but are not reported in this thesis. Two
runs of this revision had not completed at submission and contribute no number: the two arms of the distillation
ablation, submitted on 12 November.

% T13 Runs, jobs and checkpoints (Appendix A). Owner: Appendix writer (first version by the skeleton writer).
% Sources: REWRITE_PLAN.md section 4 (T13); deck explanation slide 30 (D slide 30); handover section 8 "Job index",
% 0a (state on 18 Oct, NAVER baseline, model cost, partial mask, student), 0b, 3.1-3.4; data/CLUSTER_ANSWERS.md
% (jobs submitted on the evening of 18 Oct: 3615018, 3615019, 3615020_[2,4]; 3615020_[2,4] returned 22:40;
% partial mask 3614042 "19 Oct, 00:30", its Relative Human job 3616821 "19 Oct, 01:20"; Student 3614191 done
% "19 Oct morning", log logs/exp8_distill_vitb_3614191.out).
% Trimmed 19 Oct to the jobs whose numbers, tables or figures appear in the thesis (each id below is cited in a
% numbers_*.tex comment, a table or a chapter). Dropped rows and ids (recover with
% git show 76aa1a4:docs/thesis/tables/T13_runs.tex): Mixed-L epoch-149 trace 3571964, re-fit test 3589270 and
% crashes 3587852 / 3587853, Relative Human smoke test 3608336; the "What" column merged into the first.
% The only place in the thesis where job ids appear (plan section 5), together with T19 and Appendix A's text.
% Paths are relative to the project folder on the cluster (/netscratch/najib/multi-hmr). Log = logs/<name>_<job>.out;
% array jobs log as <name>_<job>_<task>.out. Exception: the video and image launchers log as
% logs/video-compare-<job>.out and logs/image-compare-<job>.out (run_video_compare.sh:13, run_image_compare.sh:13).
% Added 6 Nov: Student 60 epochs (3623342; CA "6 Nov: 60-epoch student final table"); loss curves (login node, no job;
% scripts/analysis/plot_loss_curves.py docstring; CA "6 Nov: loss curves"); temperature sweep (exp9_kd_temperature.sh
% header "OUTPUTS"; scripts/analysis/kd_temperature_sweep.py docstring; job ids pending, open question F8); video job
% 3642415 (numbers_video.tex; Appendix A, sec:app-video); five-panel video (slurm/demo/run_video_compare.sh:98-100,
% docs/video_inference.md:36-60 and its output table; job id pending, open question F6); image comparison
% (slurm/demo/run_image_compare.sh:81, docs/image_comparison.md "Output"; job id and photo pending, open question F7);
% thesis copies F23/F25/F26 (artefacts of 6 Nov: F25 frame 178 and F26 cut from the final deck, slides 18 and 11).
% 12 Nov: caption shortened (path key moved to sec:app-runs); "Student, 60 epochs" -> Student-60; sweep row updated
% (CA "8 Nov: exp9 temperature sweep done": T=1,2,4,8 completed, T=16 3717598 trained but failed after training on an
% HTTP 504 time-out); rows added for the distillation ablation (two arms; submission not confirmed, no job ids) and the
% student's Relative Human evaluation (pending), brief section 5. Fix pass 12 Nov (D1): ablation rows "not completed at
% submission" (launcher slurm/experiments/exp10_kd_ablation.sh, record file), Relative Human row of the student and the
% sweep's pending-validation clause removed; V1 -> holdout protocol.
% 13 Nov: ablation rows completed (CA "13 Nov: exp10 results" A: jobs 3744297, 3744307, submitted 12 Nov 22:50, git
% 7b7edf1, record logs/exp10_kd_ablation_20261112_225014.txt); sweep validation table added to the sweep row (CA 13 Nov B).
\begingroup
\expandafter\def\expandafter\UrlBreaks\expandafter{\UrlBreaks\do\_\do\-}% let \path break after _ and - here
\scriptsize
\setlength{\tabcolsep}{2.5pt}
\begin{longtable}{@{}>{\raggedright}p{3.75cm}>{\raggedright}p{2.35cm}>{\raggedright\arraybackslash}p{6.5cm}@{}}
\caption[Runs, jobs, logs and checkpoints]{Slurm jobs behind the numbers, tables and figures of the thesis, with
log name and checkpoint or output. Run identifiers as in Table~\ref{tab:names}; path conventions in
Section~\ref{sec:app-runs}.}
\label{tab:runs}\\
\toprule
Model or item: what the job did & Job(s) & Log name; checkpoint or output \\
\midrule
\endfirsthead
\toprule
Model or item: what the job did & Job(s) & Log name; checkpoint or output \\
\midrule
\endhead
\bottomrule
\endfoot
\multicolumn{3}{@{}l}{\emph{Models}} \\
Base: holdout trace (PVE in the holdout protocol, placement) & 3570222; re-trace 3571596 & \path{diag_holdout}; \path{logs/anny_model/anny_s2_shape_v5/checkpoints/00099.pt} \\
Mixed-L: holdout trace & 3571963 & \path{diag_holdout}; \textsf{M}\path{/anny_teacher_v9_mixed_00062.pt} \\
Fresh-heads-L: holdout trace & 3571794 & \path{diag_holdout}; \path{logs/anny_model/anny_v10_scratch_mixed/checkpoints/00290.pt} \\
Final: training, holdout trace, 3DPW test & 3611336 (smoke test 3611333) & \path{exp7_final_model}; \textsf{E7}\path{/v5_pc_dist5_bonefix_v23_e10/train/checkpoints/00009.pt}, copy \textsf{M}\path{/exp7_v5_final_00009.pt} \\
Control: the same & 3611339 & \path{exp7_final_model}; \textsf{E7}\path{/v5_root_dist1_nofix_v23_e10/} \\
Final-masked: the same & 3612325 & \path{exp7_final_model}; \textsf{E7}\path{/v5_pc_dist5_bonefix_v23_e10_shpmask/} \\
Final-partial-mask: the same (done 19 Oct) & 3614042 & \path{exp7_final_model} (holdout table lines 191--215); \textsf{E7}\path{/v5_pc_dist5_bonefix_v23_e10_shpmask_noage/} \\
Student: distillation from Final, 30 epochs, own evaluation (done 19 Oct) & 3614191 (smoke test 3613926) & \path{exp8_distill_vitb}; \path{results/experiments/exp8_distill_vitb/student_vitb_from_v5_pc_dist5_bonefix_v23_e10_00009_e30/train/checkpoints/00029.pt}; \path{final_table.json} in the same run folder \\
Multi-HMR-Anny: holdout and 3DPW test with the evaluation code of this thesis & 3614192 (3613928 crashed: FFN width) & \path{eval-test3dpw} (task 0); \textsf{M}\path{/multiHMR_672_L_anny.pt} \\
Model cost: parameters, GFLOPs, latency of Multi-HMR-Anny and Final (A100) & 3614947 & \path{model_cost}; \path{results/diagnostics/model_cost_naver_vs_final.json} \\
\addlinespace
\multicolumn{3}{@{}l}{\emph{Training length, distillation temperature and distillation ablation}} \\
Student-60: the same training configuration trained twice as long, own evaluation & 3623342 & \path{exp8_distill_vitb}; \path{results/experiments/exp8_distill_vitb/student_vitb_from_v5_pc_dist5_bonefix_v23_e10_00009_e60/train/checkpoints/00059.pt}; \path{final_table.json} in the same run folder \\
Loss curves of the 30- and 60-epoch runs, from their tensorboard files (Figure~\ref{fig:loss-curves}) & none (login node) & \path{scripts/analysis/plot_loss_curves.py}; \path{results/experiments/exp9_kd_temperature/loss_curves/<run>_loss_curves.png}, \path{.pdf}, \path{.csv}; thesis figure drawn from the CSV copies \path{data/loss_curves_student_e30.csv}, \path{_e60.csv} by \path{figures/scripts/mk_f23_loss_curves.py} \\
Temperature sweep, \KDSweepArms{} runs (submitted 6 Nov; $T = 1, 2, 4, 8$ finished; $T = 16$ trained to the end, but its post-training evaluation failed on a network time-out) & 3717507, 3717530, 3717538, 3717581, 3717598 ($T = 1, 2, 4, 8, 16$) & \path{exp8_distill_vitb} (one job per temperature, submitted by \path{slurm/experiments/exp9_kd_temperature.sh}; record \path{logs/exp9_kd_temperature_<date>.txt}); collector \path{scripts/analysis/kd_temperature_sweep.py} writes \path{results/experiments/exp9_kd_temperature/kd_temperature_sweep.csv}, \path{.png}, \path{.pdf}; final-epoch validation values of all five runs (Table~\ref{tab:sweep-validation}) from \path{scripts/analysis/sweep_validation_table.py}, \path{results/experiments/exp10_kd_ablation/validation/sweep_validation_table.csv}, thesis copy \path{data/sweep_validation_table.csv}\\
Distillation ablation, the student without distillation losses (Student-noKD): both distillation losses off, teacher heads copied (submitted 12 Nov, 22:50, git 7b7edf1; completed) & 3744297 & \path{exp8_distill_vitb}, submitted by \path{slurm/experiments/exp10_kd_ablation.sh} (record \path{logs/exp10_kd_ablation_20261112_225014.txt}); \path{results/experiments/exp8_distill_vitb/student_vitb_from_v5_pc_dist5_bonefix_v23_e10_00009_e30_nokd/}: \path{final_table.json}, \path{bone_fix.json}; collector \path{scripts/analysis/kd_ablation_table.py} writes \path{results/experiments/exp10_kd_ablation/kd_ablation_table.csv}; thesis copy \path{data/kd_ablation_table.csv} \\
Distillation ablation, the ViT-B model without teacher (ViT-B-noTeacher): same data and schedule, no teacher (submitted 12 Nov, 22:50, git 7b7edf1; completed) & 3744307 & as the row above, run folder \path{..._e30_supervised/} \\
\addlinespace
\multicolumn{3}{@{}l}{\emph{Evaluations and demos}} \\
3DPW test (v23): Base, Mixed-L, neck pilot (Final and Control in their own jobs) & 3612073, 3612074, 3612075 & \path{eval-test3dpw} (task 0), lines 117--125 \\
3DPW test (v23): Student-S/6c, Student-S/final, also AnnyOne test-100 & 3615020\_2, 3615020\_4 & \path{eval-test3dpw}, line 119 (\_2) and 118 (\_4) \\
Relative Human: Base, neck pilot, Final, Final-masked, Final-partial-mask & 3611285\_0, \_1, 3612131, 3613616, 3616821 & \path{eval_rh}; \path{results/diagnostics/relative_human_val_<tag>/} \\
Demos: neck pilot, Final, Final-masked, each next to Base & 3581312, 3612078, 3613617 & \path{demo-compare}; \path{demo_montage_photos/}, \path{demo_montage_3dpwtest/} \\
Demos: predicted fingers (Base, Final); rest pose (Base, Final, Final-masked) & 3615018, 3615019 & \path{demo-compare}; scripts \path{slurm/demo/run_demo_compare_hands.sh}, \path{_rest.sh} (copies made on the cluster) \\
Video, raw Student and Multi-HMR-Anny on \path{man_walking.mp4} (frame and match counts) & 3642415 & \path{video-compare} (log \path{logs/video-compare-3642415.out}); \path{results/video_compare/video_3642415/comparison.mp4} \\
Video, five panels with temporal filter and landmark alignment (Figure~\ref{fig:video-panels}) & \VideoJob & \path{video-compare}; \path{slurm/demo/run_video_compare.sh} with \path{TEMPORAL=ema}, \path{ALIGNMENT=landmarks} and \path{LANDMARKS=<json>}; landmarks from \path{scripts/video/extract_pose_landmarks.py} (separate MediaPipe 0.10.14 environment); \path{results/video_compare/video_3657550/}: \path{comparison.mp4}, \path{timing.csv}, \path{student_temporal/}, \path{student_aligned/alignment.json}; thesis copy \path{figures/img/F25_video_panels.jpg} (frame \VideoFigureFrame{} of the video in the final presentation) \\
Stage photo, Student and Multi-HMR-Anny (Figure~\ref{fig:image-compare}) & 3657637 & \path{image-compare}; \path{slurm/demo/run_image_compare.sh}, which runs \path{demo_image_compare.py}; \path{results/image_compare/images_3657637/comparison/}; thesis copy \path{figures/img/F26_image_compare.jpg} (cut from the final presentation) \\
\addlinespace
\multicolumn{3}{@{}l}{\emph{3DPW labels}} \\
Label check before and after the re-fit & 3575369, 3611284 (v23) & \path{check_3dpw_conversion}; \path{results/diagnostics/check_3dpw_conversion_v23/} \\
Re-fit with converter v2.3, 87 person-sequences & 3590385, 3602813 & \path{smpl2anny-all}; \path{threedpw_anny_v23/labels/}; \path{logs/check_refit_all_v23.log} \\
Label rebuilds from the old fits & 3579266, 3581320 & \path{3dpw-labels}; \path{threedpw_anny_full/labels_rebuilt/} \\
\addlinespace
\multicolumn{3}{@{}l}{\emph{Pilots and diagnostics (outputs in \path{results/experiments/} and \path{results/diagnostics/})}} \\
Phenotypes; youngest 50 training people & 3569948; 3569950 & \path{diag_pheno}; \path{diag_young} \\
Real-photo ages with the crop camera: Base, sampler pilot, its control & 3569949, 3572419, 3572420 & \path{diag_cropK} \\
Depth pilots, 2 epochs ($\times$1, control, $\times$5, $\times$20) and 10 epochs ($\times$1, control, $\times$5) & 3571598, 3571599, 3572418, 3572483; 3571933, 3571934, 3572482 & \path{exp1_depth_anchor} \\
Sampler pilot and control; batch ages only & 3571955, 3571956; 3571601 & \path{exp2_children} \\
Head-loss pilot; two identical controls (noise) & 3571957; 3571958, 3571960 & \path{exp3_neck_global_rot}; \path{exp4_fingers} (3571960) \\
Finger pilot (control 3571960: crop trace, hand size) & 3571959 & \path{exp4_fingers} \\
Geodesic pilot and its L1 control & 3573520, 3573521 & \path{exp5_neck_fix} \\
Per-bone dumps of eight checkpoints & 3573518, 3573519 & \path{diag_neck_bones}; \path{results/diagnostics/neck_bones/} \\
Bone correction: Base (with the neck pilot), Fresh-heads-L & 3579014, 3579015 & \path{exp6_bone_fix}; \path{results/experiments/exp6_bone_fix/anny_s2_shape_v5_00099_n200_s0/bone_fix.json}; neck pilot \textsf{M}\path{/exp6_v5_bonefix_00001.pt} \\
\end{longtable}
\endgroup


The scripts behind these jobs and the exact commands that produced the thesis numbers are listed in the file
\path{docs/thesis/REPRODUCE.md} of the repository.\label{sec:app-scripts}\label{sec:app-commands} Each command runs
from the login node in the repository clone. The flags each training run passes are in Appendix~\ref{app:flags}.

\section{What cannot be reproduced exactly}
\label{sec:app-repro-limits}

The training code sets no random seed. The \texttt{-{}-seed} of the holdout trace only selects the training images on
which the bone correction is estimated. The student's script passes its seed to that trace and nowhere else
(\path{exp8_distill_vitb.sh}, lines 257 and 1142). The sweep runs through the same script, so this also holds for the
60-epoch student and for every temperature. The order in which training images are drawn (with replacement)
therefore differs between two runs of the same command, as does the number of optimiser steps skipped for
non-finite gradients (\MethodFinalNaNSteps{} in Final, 10 in Control, out of $\num{\MethodFinalSteps}$). Two identical
2-epoch pilots differed by \NoisePVE\,mm PVE, \NoiseHeadDeg\,deg head direction and \DiagNoiseRecall{} recall points
(Table~\ref{tab:pilot-noise}); they do not establish an interval for a repeat of Final or Control
(Section~\ref{sec:method-variation}).

Final and Control also ran on different GPU types (Final on an A100, Control on an A40, both with TF32). This changes
the speed, and it can change the last digits of the arithmetic. Evaluating a fixed checkpoint involves no sampling;
whether it repeats exactly across GPU types was not tested.

Apart from the temperature, the sweep runs evaluate after every epoch instead of every fifth and also compute the
training loss on the validation images (\texttt{-{}-val\_loss 1}). Both additions only read values out of the model
and change no weight. The sweep's own $T=4$ run therefore repeats the reported training configuration and gives a second estimate of
run-to-run variation (Section~\ref{sec:method-variation}).

Three further items cannot be repeated. The history scripts of the L line and of Student-S were edited after their
runs, so their training configurations in Table~\ref{tab:names} are as recorded there, not re-run. \path{train.py} keeps only the
last two or three checkpoints of a run, so intermediate epochs are lost. The checkpoints, the AnnyOne and 3DPW data
and the Relative Human annotations are not in the repository, and 3DPW and Relative Human come with their own
licences.

\section{Records of the video comparison}
\label{sec:app-video}

The video comparison of job 3642415 is documented by the log \path{video-compare-3642415.out} and the matching error
log. It used the 30-epoch checkpoint of the student (Student)
\path{results/experiments/exp8_distill_vitb/student_vitb_from_v5_pc_dist5_bonefix_v23_e10_00009_e30/train/checkpoints/00029.pt}
with the correction selected by that run's \path{student_fix_arg.txt}. The released baseline (Multi-HMR-Anny) was
\path{models/multiHMR/multiHMR_672_L_anny.pt} and the input was \path{man_walking.mp4}. Both checkpoints loaded with
100\,\% state coverage, meaning every saved weight found its place in the model. The renderer reported an NVIDIA L40S
and passed its test mesh.

The final log lines report \VideoFrames{} processed frames for each model, \VideoStudentMatched{} student matches and
\VideoBaselineMatched{} baseline matches. The output is \path{results/video_compare/video_3642415/comparison.mp4}. The
job ran before the temporal filter and the later timing code existed. Its log therefore gives neither the inference
speed of the two models on the same frames nor any measure of frame-to-frame change in pose or body shape.

The five-panel video of Section~\ref{sec:res-video} came from a later run of the same launcher with
\path{TEMPORAL=ema} and \path{ALIGNMENT=landmarks}. It adds the filtered and the aligned student panels
(Section~\ref{sec:method-temporal}) and a frame-rate label on every panel. This run is job \VideoJob{} (\path{results/video_compare/video_3657550/}). It ran
on one \VideoGPU{} in fp32 with TF32 at 672 pixels, timed \VideoMeasuredFrames{} frames after \VideoWarmup{}
warm-up frames, and used the filter constants of Section~\ref{sec:method-temporal}. Its own summaries give the same
matched-frame counts as job 3642415: \VideoStudentMatched{} of \VideoFrames{} for the student and
\VideoBaselineMatched{} for the baseline.

The frame-rate label on each panel counts a different amount of work (Table~\ref{tab:video-fps}; definitions in
\path{docs/video_inference.md}, lines 387--440). \path{timing.csv} has one row per model panel, and the input panel's
playback rate is a column of those rows. The raw student and the
baseline show \texttt{inference\_fps}, model inference alone. The filtered panel shows \texttt{processing\_fps} from
\path{student_temporal/summary.json}, which also counts the filter and the Anny re-decode. The aligned panel shows the
same quantity from \path{student_aligned/summary.json}, which counts the filter work as well as the alignment fit, but not the
offline landmark extraction.

The video implementation and the EGL correction are in commits \texttt{1d5621b} and \texttt{7b0b5e1}, and the revised
FPS implementation is \texttt{a0fa59f}. The video comparison is launched through \path{slurm/demo/run_video_compare.sh}, and
\path{docs/video_inference.md} describes how to run it, where each timing starts and stops, and which files it
writes. The stage photo of Section~\ref{sec:res-image-compare} was made by \path{demo_image_compare.py}, normally
launched through \path{slurm/demo/run_image_compare.sh}, whose defaults are the same student, correction and baseline
(\path{docs/image_comparison.md}). It is job 3657637, with its output in
\path{results/image_compare/images_3657637/}.

\endgroup
